% ─────────────────────────────────────────────────────────────────────────────
%  estilo.tex — O DESIGN DO REINO, um só para os três documentos.
%
%  Era o enredo que o tinha inteiro (fancyhdr, titlesec, as sete fontes gk*, as
%  cores dos personagens); a teoria e o catálogo tinham metade cada um, e a
%  teoria nem sequer carregava o tcolorbox que usava 8 vezes — as opções da
%  caixa saíam IMPRESSAS no PDF publicado. Unificar corrige isso de graça.
%
%  Todos os documentos são `report`: o main.tex junta os três, e cada um
%  continua a compilar sozinho via `subfiles`.
%
%  REGRA: aqui tudo é \providecommand. Um \newcommand repetido é erro fatal
%  quando os três corpos entram no mesmo documento.
% ─────────────────────────────────────────────────────────────────────────────

\usepackage[utf8]{inputenc}\usepackage[T1]{fontenc}\usepackage[brazil]{babel}
\usepackage{amsmath,amssymb,amsthm,mathtools}
\usepackage{geometry}\usepackage{xcolor}\usepackage{booktabs}
\usepackage{array}\usepackage{longtable}
% mdframed: caixa com moldura que QUEBRA entre páginas. O \fbox{minipage} do enunciado
% não quebrava — e pior, anulava o longtable lá dentro, que dentro de minipage não parte.
\usepackage{mdframed}\usepackage{textcomp}

% ── o sumário passou dos 99 capítulos: 3 dígitos não cabiam e o título subia
%    por cima do número. tocloft vem ANTES do hyperref.
\usepackage{tocloft}
\setlength{\cftchapnumwidth}{3.2em}
\setlength{\cftsecnumwidth}{4.0em}
\setlength{\cftsubsecnumwidth}{5.0em}
\setlength{\cftpartnumwidth}{3.2em}
\setlength{\cftsecindent}{3.2em}
\setlength{\cftsubsecindent}{7.2em}
\makeatletter\renewcommand{\@pnumwidth}{2.4em}\renewcommand{\@tocrmarg}{3.4em}\makeatother
\usepackage{hyperref}
\hypersetup{colorlinks=true, linkcolor=black!70, urlcolor=black!70}

\usepackage{fvextra}\usepackage[htt]{hyphenat}\usepackage[expansion=false]{microtype}
\usepackage{ragged2e}\usepackage{fancyhdr}\usepackage{titlesec}
\usepackage{tcolorbox}\tcbuselibrary{breakable,skins}
\usepackage{fancyvrb}\usepackage{truncate}\usepackage{xspace}
\usepackage{tikz}\usepackage{graphicx}
\usepackage{skak}   % o enredo compõe tabuleiros (\fenboard, \showboard) — é um livro de xadrez
\hyphenpenalty=50 \tolerance=2000
\sloppy \emergencystretch=3em

% ── os caracteres Unicode literais que os corpos têm em modo texto ──
\DeclareUnicodeCharacter{2081}{\ensuremath{_1}}
\DeclareUnicodeCharacter{2082}{\ensuremath{_2}}
\DeclareUnicodeCharacter{2083}{\ensuremath{_3}}
\DeclareUnicodeCharacter{2212}{\ensuremath{-}}
\DeclareUnicodeCharacter{2229}{\ensuremath{\cap}}
\DeclareUnicodeCharacter{0394}{\ensuremath{\Delta}}
\DeclareUnicodeCharacter{22C8}{ com }

% ── a semente da estrela: respiro (dobra/3), ascendente 17/20, descendente 1/4, traço 400 ──
% Só as sementes são fixas: todo o resto da tipografia DERIVA delas pela espiral.
% O pdflatex ignora (comando vazio); o tradutor lê-a como config (le_semente).
\providecommand{\gksemente}[5]{}
\gksemente{3}{17}{20}{4}{400}

% ── as sete fontes e as quatro cores da casa ──
\providecommand{\gknota}{\fontsize{7.62}{11.05}\selectfont}
\providecommand{\gkcod}{\fontsize{8.94}{12.97}\selectfont}
\providecommand{\gktexto}{\fontsize{10.50}{15.22}\selectfont}
\providecommand{\gksub}{\fontsize{12.33}{17.87}\selectfont}
\providecommand{\gksec}{\fontsize{14.47}{20.98}\selectfont}
\providecommand{\gkcap}{\fontsize{16.99}{24.63}\selectfont}
% o degrau k=6 não é escolha: o relógio de 8 marcas fá-lo o DUAL do k=2 (gksub), e sem ele
% a involução da escala não fecha — `x+x†` dá ímpar nos sete e a divisão por dois falha.
% Corpo 19,95 = 7,62·φ^(6/3), entrelinha 28,92 pela mesma razão 1,4497 de todos os outros.
\providecommand{\gkparte}{\fontsize{19.95}{28.92}\selectfont}
\providecommand{\gktit}{\fontsize{23.42}{33.95}\selectfont}
\definecolor{ouro}{HTML}{B8912F}\definecolor{ouroclaro}{HTML}{F3E9CE}
\definecolor{tinta}{HTML}{1E1B16}\definecolor{regua}{HTML}{6E6858}

\fvset{breaklines=true, breakanywhere=true, fontsize=\gkcod,
       frame=leftline, framerule=1.2pt, rulecolor=\color{ouro}, framesep=3mm, xleftmargin=4mm}
\hyphenation{PON-TRY-A-GIN Pon-try-a-gin CLIF-FORD Clif-ford ELE-TRO-MAG-NE-TI-CO
             ele-tro-mag-ne-ti-co Pri-mei-ri-da-de Se-cun-di-da-de Ter-cei-ri-da-de
             es-pa-co-tem-po-ral te-les-co-pi-co de-fle-xi-vo cris-ta-li-no
             MYHILL-NE-RODE ARTIN-WHA-PLES POIN-SE-TI-CA mo-du-la-cao au-to-si-mi-lar}

% ── os títulos dourados e o cabeçalho ──
\titleformat{\chapter}[display]
  {\normalfont\bfseries\color{tinta}}
  {\color{ouro}\gknota\scshape\chaptertitlename\ \thechapter}{6pt}
  {\gktit}[{\vspace{2mm}\color{ouro}\titlerule[1.2pt]}]
\titlespacing*{\chapter}{0pt}{-14pt}{22pt}
\titleformat{\section}{\normalfont\gksec\bfseries\color{tinta}}{\color{ouro}\thesection}{0.7em}{}
\titleformat{\subsection}{\normalfont\gksub\bfseries\color{regua}}{\color{ouro}\thesubsection}{0.6em}{}

\pagestyle{fancy}\fancyhf{}
\renewcommand{\chaptermark}[1]{\markboth{\thechapter.\ #1}{}}
\newlength{\gkwmarca}\newlength{\gkwcap}
\setlength{\gkwmarca}{0.24\textwidth}
\setlength{\gkwcap}{0.74\textwidth}
\fancyhead[L]{\parbox[b]{\gkwmarca}{\gknota\scshape\color{regua}Reino Dourado}}
\fancyhead[R]{\parbox[b]{\gkwcap}{\raggedleft\gknota\scshape\color{regua}%
  \truncate{\gkwcap}{\nouppercase{\leftmark}}}}
\fancyfoot[C]{\gknota\color{regua}\thepage}
\renewcommand{\headrulewidth}{0.4pt}
\renewcommand{\headrule}{\color{ouro}\hrule height 0.4pt}
\setlength{\headheight}{16pt}
\addtolength{\topmargin}{-4pt}
\renewcommand{\arraystretch}{1.25}\setlength{\tabcolsep}{4pt}
\geometry{margin=2.6cm}

% ── os teoremas: UM contador para todos, preso ao capítulo ──
%    A teoria numerava por \section e o catálogo tinha `obs` num contador à
%    parte. No documento único isso davam duas séries a competir.
\theoremstyle{plain}
\newtheorem{teorema}{Teorema}[chapter]
\newtheorem{proposicao}[teorema]{Proposição}
\newtheorem{corolario}[teorema]{Corolário}
\newtheorem{principio}[teorema]{Princípio}
\theoremstyle{definition}
\newtheorem{definicao}[teorema]{Definição}
\newtheorem{defn}[teorema]{Definição}
\newtheorem{exemplo}[teorema]{Exemplo}
\newtheorem{conjectura}[teorema]{Conjectura}
\theoremstyle{remark}
\newtheorem{obs}[teorema]{Observação}
\newtheorem{observacao}[teorema]{Observação}
\renewcommand{\proofname}{Demonstração}

% ── as macros matemáticas, de todos os três ──
\providecommand{\code}[1]{\texttt{\footnotesize #1}}
\providecommand{\R}{\mathbb{R}}\providecommand{\C}{\mathbb{C}}
\providecommand{\Q}{\mathbb{Q}}\providecommand{\Z}{\mathbb{Z}}
\providecommand{\N}{\mathbb{N}}\providecommand{\F}{\mathbb{F}}
\providecommand{\Hq}{\mathbb{H}}\providecommand{\Hyp}{\mathbb{H}}
\providecommand{\Oo}{\mathbb{O}}\providecommand{\Pj}{\mathbb{P}^1}
\providecommand{\GL}{\mathrm{GL}}\providecommand{\GF}{\mathrm{GF}}
\providecommand{\Gal}{\mathrm{Gal}}\providecommand{\PSL}{\mathrm{PSL}}
\providecommand{\SL}{\mathrm{SL}}\providecommand{\Sig}{\Sigma}
\providecommand{\im}{\operatorname{Im}}
\providecommand{\sen}{\operatorname{sen}}\providecommand{\senh}{\operatorname{senh}}
\providecommand{\dm}{\operatorname{dm}}\providecommand{\tr}{\operatorname{tr}}
\providecommand{\disc}{\operatorname{disc}}\providecommand{\Pf}{\operatorname{Pf}}
\providecommand{\cf}[2]{[#1;#2]}
\providecommand{\Comp}[2]{C_{#1,#2}}
\providecommand{\Mat}[1]{\begin{psmallmatrix}#1&1\\1&0\end{psmallmatrix}}
\providecommand{\conv}{\circledast}
\providecommand{\Tiffany}{\textsc{Tiffany}}
\providecommand{\BAI}{\ensuremath{\mathrm{BAI}}}
\providecommand{\CalC}{\ensuremath{\mathcal{C}}}
\providecommand{\Collapse}{\ensuremath{\mathrm{Collapse}}}
\providecommand{\dep}{\ensuremath{\mathrm{dep}}}
\providecommand{\Fp}{\ensuremath{\mathbb F}}
\providecommand{\est}{\ensuremath{\mathrm{est}}}
% \FF pertence ao skak (macro interna dele). O 𝔽 do catálogo chama-se \Fbb.
\providecommand{\Fbb}{\ensuremath{\mathbb{F}}}\providecommand{\GG}{\ensuremath{\mathbb{G}}}
\providecommand{\RR}{\ensuremath{\mathbb{R}}}\providecommand{\ZZ}{\ensuremath{\mathbb{Z}}}
\providecommand{\citado}[1]{\emph{#1}}
\providecommand{\medido}[1]{\textbf{#1}}
\providecommand{\SI}[2]{#1\,#2}

% ── as cores dos personagens, e os seus nomes ──
\definecolor{cYasmin}{HTML}{EC0B8B}\definecolor{cBenjamim}{HTML}{E8890C}
\definecolor{cDark}{HTML}{6A1FD0}\definecolor{cJoaquim}{HTML}{0AA6C2}
\definecolor{cVenom}{HTML}{9D00C6}\definecolor{cLock}{HTML}{00A878}
\definecolor{cOperador}{HTML}{B8860B}\definecolor{cRosa}{HTML}{FF4FA3}
\definecolor{cFita}{HTML}{A98F5B}\definecolor{cMonte}{HTML}{C41E3A}
\definecolor{cVale}{HTML}{1B6CA8}\definecolor{cOuroDorme}{HTML}{8A6D0B}
\definecolor{cQuina}{HTML}{6B8E23}\definecolor{cLua}{HTML}{7E8AA2}
\providecommand{\Yas}{{\color{cYasmin}\large Yasmin}\xspace}
\providecommand{\Ben}{{\color{cBenjamim}\large Benjamim}\xspace}
\providecommand{\Drk}{{\color{cDark}\large Dark}\xspace}
\providecommand{\Joa}{{\color{cJoaquim}Joaquim}\xspace}
\providecommand{\Ven}{{\color{cVenom}\large Venom}\xspace}
\providecommand{\Lck}{{\color{cLock}Lock}\xspace}
\providecommand{\Opr}{{\color{cOperador}\small Operador}\xspace}
\providecommand{\Fit}{{\color{cFita}\large Fita}\xspace}
\providecommand{\Mon}{{\color{cMonte}\large Monte}\xspace}
\providecommand{\Val}{{\color{cVale}\large Vale}\xspace}
\providecommand{\Odo}{{\color{cOuroDorme}\large Ouro-que-Dorme}\xspace}
\providecommand{\Qui}{{\color{cQuina}\large Quina}\xspace}
\providecommand{\Lua}{{\color{cLua}\large Lua}\xspace}
\providecommand{\hino}[2]{\begin{flushright}\small\itshape\color{black!55}
sobre \textbf{#1}\\ #2\end{flushright}\medskip}

% ── o efeito de ONDA: a melodia = seno, a amplitude = áurea, a cor cicla nos 7 ──
\providecommand{\ondacor}[1]{\ifcase#1 cYasmin\or cJoaquim\or cVenom\or cBenjamim\or cLock\or cDark\or cOperador\fi}
\providecommand{\onda}[1]{\leavevmode\foreach [count=\i] \s in {#1}{%
 \pgfmathsetmacro\amp{2.7*(0.618)^(int((\i-1)/7))}%
 \pgfmathsetmacro\hh{\amp*sin(\i*47)}%
 \pgfmathtruncatemacro\ci{int(mod(\i-1,7))}%
 \raisebox{\hh pt}{\color{\ondacor{\ci}}\s}\thinspace}}

% ── A CAIXA DA CASA: o que os três usam para o que fica ──
%    \gkcaixa{...} é a forma curta; o tcolorbox cru continua a valer.
\providecommand{\gkcaixa}[1]{\begin{tcolorbox}[colback=ouroclaro!35,colframe=ouro,
  boxrule=0pt,leftrule=2pt,arc=0pt,left=4mm,right=3mm,top=2mm,bottom=2mm,breakable]%
  \itshape #1\end{tcolorbox}}

% ── A CAPA, uma só, parametrizada pelo documento ────────────────────────────
%    Régua de largura FIXA (15.4cm), não \textwidth.
\providecommand{\gkcapa}[3]{%
\providecommand{\Prj}{\mathbb{P}}
\providecommand{\Trf}{\mathcal{T}}
  \title{\vspace{-0.6cm}
  {\color{ouro}\rule{15.4cm}{1.4pt}}\\[6mm]
  \textbf{\gktit\color{tinta} Reino Dourado}\\[3mm]{\gksec\itshape\color{regua} Golden Kingdom}\\[8mm]
  {\gkcap\scshape\color{ouro} #1}\\[5mm]
  {\gksub\color{regua} #2}\\[2mm]
  {\gktexto\color{regua}\itshape #3}\\[7mm]
  {\color{ouro}\rule{15.4cm}{0.6pt}}\\[9mm]{\gknota\color{regua}\begin{minipage}{15.4cm}\setlength{\parindent}{0pt}%
  \textbf{Aviso de Isenção Algorítmica:} O universo, as leis e as entidades contidas nestas páginas ---
  incluindo felinos matriciais, aves de rapina algébricas e amuletos de controle de fluxo --- são
  construções de pura ficção formal. Esta obra não descreve a realidade física, não serve como
  engenharia reversa do cosmos e não constitui doutrina religiosa. Qualquer semelhança com bugs reais do seu
  sistema operacional é mera coincidência (ou falta de reza). Prossiga por sua própria conta e
  risco.\end{minipage}}}
  \author{}\date{}}

%% mdc/lcm como operadores (usados na Teoria e no Catálogo)
\DeclareMathOperator{\lcm}{lcm}
